\documentclass[10pt,a4paper]{amsart}
\usepackage{amsmath,amssymb,amsthm,graphicx}
% For commutative diagrams we use Xy-pic
\usepackage[all]{xy}

% We use 2cell for 2-commutative diagrams.
\xyoption{2cell}
% Theorem environments.
%
\theoremstyle{plain}
\newtheorem{theorem}[subsection]{Theorem}
\newtheorem{proposition}[subsection]{Proposition}
\newtheorem{lemma}[subsection]{Lemma}

\theoremstyle{definition}
\newtheorem{definition}[subsection]{Definition}
\newtheorem{example}[subsection]{Example}
\newtheorem{exercise}[subsection]{Exercise}
\newtheorem{situation}[subsection]{Situation}

\theoremstyle{remark}
\newtheorem{remark}[subsection]{Remark}
\newtheorem{remarks}[subsection]{Remarks}

\numberwithin{equation}{subsection}

% Macros
%
\def\lim{\mathop{\mathrm{lim}}\nolimits}
\def\colim{\mathop{\mathrm{colim}}\nolimits}
\def\Spec{\mathop{\mathrm{Spec}}}
\def\Hom{\mathop{\mathrm{Hom}}\nolimits}
\def\Ext{\mathop{\mathrm{Ext}}\nolimits}
\def\SheafHom{\mathop{\mathcal{H}\!\mathit{om}}\nolimits}
\def\SheafExt{\mathop{\mathcal{E}\!\mathit{xt}}\nolimits}
\def\Sch{\mathit{Sch}}
\def\Mor{\mathop{\mathrm{Mor}}\nolimits}
\def\Ob{\mathop{\mathrm{Ob}}\nolimits}
\def\Sh{\mathop{\mathit{Sh}}\nolimits}
\def\NL{\mathop{N\!L}\nolimits}
\def\CH{\mathop{\mathrm{CH}}\nolimits}
\def\proetale{{pro\text{-}\acute{e}tale}}
\def\etale{{\acute{e}tale}}
\def\QCoh{\mathit{QCoh}}
\def\Ker{\mathop{\mathrm{Ker}}}
\def\Im{\mathop{\mathrm{Im}}}
\def\Coker{\mathop{\mathrm{Coker}}}
\def\Coim{\mathop{\mathrm{Coim}}}

% Boxtimes
%
\DeclareMathSymbol{\boxtimes}{\mathbin}{AMSa}{"02}

%
% Macros for moduli stacks/spaces
%
\def\QCohstack{\mathcal{QC}\!\mathit{oh}}
\def\Cohstack{\mathcal{C}\!\mathit{oh}}
\def\Spacesstack{\mathcal{S}\!\mathit{paces}}
\def\Quotfunctor{\mathrm{Quot}}
\def\Hilbfunctor{\mathrm{Hilb}}
\def\Curvesstack{\mathcal{C}\!\mathit{urves}}
\def\Polarizedstack{\mathcal{P}\!\mathit{olarized}}
\def\Complexesstack{\mathcal{C}\!\mathit{omplexes}}
% \Pic is the operator that assigns to X its picard group, usage \Pic(X)
% \Picardstack_{X/B} denotes the Picard stack of X over B
% \Picardfunctor_{X/B} denotes the Picard functor of X over B
\def\Pic{\mathop{\mathrm{Pic}}\nolimits}
\def\Picardstack{\mathcal{P}\!\mathit{ic}}
\def\Picardfunctor{\mathrm{Pic}}
\def\Deformationcategory{\mathcal{D}\!\mathit{ef}}
\newenvironment{reference}{}{}
\newenvironment{slogan}{}{}
\newenvironment{history}{}{}
\usepackage{cleveref}
\crefname{originaldefinition}{original definition}{original definitions}
\Crefname{originaldefinition}{Original definition}{Original definitions}
\makeatletter
\newlabel{homology-definition-filtered-complex}{{012L}{source}{}{sourcepointer-1}{}}
\newlabel{homology-definition-filtered-complex@cref}{{[originaldefinition][0][]012L}{[1][0][]source}{}{}{}}
\makeatother
\crefname{originaltheorem}{original theorem}{original theorems}
\Crefname{originaltheorem}{Original theorem}{Original theorems}
\makeatletter
\newlabel{theorem-K-injective-embedding-grothendieck}{{079P}{source}{}{sourcepointer-2}{}}
\newlabel{theorem-K-injective-embedding-grothendieck@cref}{{[originaltheorem][0][]079P}{[1][0][]source}{}{}{}}
\makeatother
\crefname{originallemma}{original lemma}{original lemmas}
\Crefname{originallemma}{Original lemma}{Original lemmas}
\makeatletter
\newlabel{derived-lemma-product-K-injective}{{0BK6}{source}{}{sourcepointer-3}{}}
\newlabel{derived-lemma-product-K-injective@cref}{{[originallemma][0][]0BK6}{[1][0][]source}{}{}{}}
\makeatother
\crefname{originallemma}{original lemma}{original lemmas}
\Crefname{originallemma}{Original lemma}{Original lemmas}
\makeatletter
\newlabel{derived-lemma-K-injective}{{070I}{source}{}{sourcepointer-4}{}}
\newlabel{derived-lemma-K-injective@cref}{{[originallemma][0][]070I}{[1][0][]source}{}{}{}}
\makeatother
\crefname{originallemma}{original lemma}{original lemmas}
\Crefname{originallemma}{Original lemma}{Original lemmas}
\makeatletter
\newlabel{derived-lemma-triangle-K-injective}{{090X}{source}{}{sourcepointer-5}{}}
\newlabel{derived-lemma-triangle-K-injective@cref}{{[originallemma][0][]090X}{[1][0][]source}{}{}{}}
\makeatother
\title{Simultaneous K-injective resolutions for a filtered complex and all its filtration subquotients}
\author{The Stacks Project Authors}
\date{stacks-project:injectives.tex:local-snapshot-20261011; GNU FDL 1.2 or later}
\begin{document}
\maketitle
\noindent\textit{Editorial note.} Public360 is bounded below and uses finite filtered-injective resolutions. This statement allows the original unbounded filtered complex and asks for simultaneousK-injectivity of all subquotients; its product/nullhomotopy construction is additional work.\par\medskip
\noindent\textit{Allowed standard prerequisites (editorial).}\begin{itemize}
\item Explicit prerequisite boundary. Grothendieck abelian categories, functorial unboundedK-injective resolutions, cones and homotopies, products ofK-injectives and split exact sequences. The local derived-product argument is retained.
\end{itemize}\medskip
\ifdefined\hypertarget\hypertarget{sourcepointer-1}{}\fi
\noindent\textit{Source reference.} Original definition 012L: focused-seeds/AGP/sources/stacks-snapshot-20261011/tags.txt, source lines 1403--1403. Stable Stacks tag for an expressly allowed foundational or notational input; the counted construction remains the two concrete core obligations. The retained local core is not replaced by this pointer. See standard\_prerequisites.\par\medskip
\ifdefined\hypertarget\hypertarget{sourcepointer-2}{}\fi
\noindent\textit{Source reference.} Original theorem 079P: focused-seeds/AGP/sources/stacks-snapshot-20261011/tags.txt, source lines 9279--9279. Stable Stacks tag for an expressly allowed foundational or notational input; the counted construction remains the two concrete core obligations. The retained local core is not replaced by this pointer. See standard\_prerequisites.\par\medskip
\ifdefined\hypertarget\hypertarget{sourcepointer-3}{}\fi
\noindent\textit{Source reference.} Original lemma 0BK6: focused-seeds/AGP/sources/stacks-snapshot-20261011/tags.txt, source lines 14549--14549. Stable Stacks tag for an expressly allowed foundational or notational input; the counted construction remains the two concrete core obligations. The retained local core is not replaced by this pointer. See standard\_prerequisites.\par\medskip
\ifdefined\hypertarget\hypertarget{sourcepointer-4}{}\fi
\noindent\textit{Source reference.} Original lemma 070I: focused-seeds/AGP/sources/stacks-snapshot-20261011/tags.txt, source lines 8957--8957. Stable Stacks tag for an expressly allowed foundational or notational input; the counted construction remains the two concrete core obligations. The retained local core is not replaced by this pointer. See standard\_prerequisites.\par\medskip
\ifdefined\hypertarget\hypertarget{sourcepointer-5}{}\fi
\noindent\textit{Source reference.} Original lemma 090X: focused-seeds/AGP/sources/stacks-snapshot-20261011/tags.txt, source lines 11425--11425. Stable Stacks tag for an expressly allowed foundational or notational input; the counted construction remains the two concrete core obligations. The retained local core is not replaced by this pointer. See standard\_prerequisites.\par\medskip
\begin{lemma}
\label{lemma-K-injective-embedding-filtration}
Let $\mathcal{A}$ be a Grothendieck abelian category.
Let $K^\bullet$ be a filtered complex of $\mathcal{A}$, see
Homology, Definition \ref{homology-definition-filtered-complex}.
Then there exists a morphism $j : K^\bullet \to J^\bullet$
of filtered complexes of $\mathcal{A}$ such that
\begin{enumerate}
\item $J^n$, $F^pJ^n$, $J^n/F^pJ^n$ and $F^pJ^n/F^{p'}J^n$ are injective
objects of $\mathcal{A}$,
\item $J^\bullet$, $F^pJ^\bullet$, $J^\bullet/F^pJ^\bullet$, and
$F^pJ^\bullet/F^{p'}J^\bullet$ are K-injective complexes,
\item $j$ induces quasi-isomorphisms
$K^\bullet \to J^\bullet$,
$F^pK^\bullet \to F^pJ^\bullet$,
$K^\bullet/F^pK^\bullet \to J^\bullet/F^pJ^\bullet$, and
$F^pK^\bullet/F^{p'}K^\bullet \to F^pJ^\bullet/F^{p'}J^\bullet$.
\end{enumerate}
\end{lemma}

\begin{lemma}
\label{lemma-derived-products}
Let $\mathcal{A}$ be a Grothendieck abelian category.
Then
\begin{enumerate}
\item $D(\mathcal{A})$ has both direct sums and products,
\item direct sums are obtained by taking termwise direct sums of
any complexes,
\item products are obtained by taking termwise products of
K-injective complexes.
\end{enumerate}
\end{lemma}

\begin{proof}
Let $K^\bullet_i$, $i \in I$ be a family of objects of $D(\mathcal{A})$
indexed by a set $I$. We claim that the termwise direct sum
$\bigoplus_{i \in I} K^\bullet_i$ is a direct sum in $D(\mathcal{A})$.
Namely, let $I^\bullet$ be a K-injective complex. Then we have
\begin{align*}
\Hom_{D(\mathcal{A})}(\bigoplus\nolimits_{i \in I} K^\bullet_i, I^\bullet)
& =
\Hom_{K(\mathcal{A})}(\bigoplus\nolimits_{i \in I} K^\bullet_i, I^\bullet) \\
& =
\prod\nolimits_{i \in I} \Hom_{K(\mathcal{A})}(K^\bullet_i, I^\bullet) \\
& =
\prod\nolimits_{i \in I} \Hom_{D(\mathcal{A})}(K^\bullet_i, I^\bullet)
\end{align*}
as desired. This is sufficient since any complex can be represented
by a K-injective complex by
Theorem \ref{theorem-K-injective-embedding-grothendieck}.
To construct the product, choose a K-injective resolution
$K_i^\bullet \to I_i^\bullet$ for each $i$. Then we claim that
$\prod_{i \in I} I_i^\bullet$ is a product in $D(\mathcal{A})$.
This follows from
Derived Categories, Lemma \ref{derived-lemma-product-K-injective}.
\end{proof}

\begin{proof}
By Theorem \ref{theorem-K-injective-embedding-grothendieck}
we obtain quasi-isomorphisms $i : K^\bullet \to I^\bullet$ and
$i^p : F^pK^\bullet \to I^{p, \bullet}$ as well as commutative diagrams
$$
\vcenter{
\xymatrix{
K^\bullet \ar[d]_i & F^pK^\bullet \ar[l] \ar[d]_{i^p} \\
I^\bullet & I^{p, \bullet} \ar[l]_{\alpha^p}
}
}
\quad\text{and}\quad
\vcenter{
\xymatrix{
F^{p'}K^\bullet \ar[d]_{i^{p'}} &
F^pK^\bullet \ar[l] \ar[d]_{i^p} \\
I^{p', \bullet} &
I^{p, \bullet} \ar[l]_{\alpha^{p p'}}
}
}
\quad\text{for }p' \leq p
$$
such that $\alpha^p \circ \alpha^{p' p} = \alpha^{p'}$
and $\alpha^{p'p''} \circ \alpha^{pp'} = \alpha^{pp''}$.
The problem is that the maps $\alpha^p : I^{p, \bullet} \to I^\bullet$
need not be injective. For each $p$ we choose an injection
$t^p : I^{p, \bullet} \to J^{p, \bullet}$ into an acyclic K-injective
complex $J^{p, \bullet}$ whose terms are injective objects of $\mathcal{A}$
(first map to the cone on the identity and then use the theorem).
Choose a map of complexes $s^p : I^\bullet \to J^{p, \bullet}$
such that the following diagram commutes
$$
\xymatrix{
K^\bullet \ar[d]_i & F^pK^\bullet \ar[l] \ar[d]_{i^p} \\
I^\bullet \ar[rd]_{s^p} & I^{p, \bullet} \ar[d]^{t^p} \\
& J^{p, \bullet}
}
$$
This is possible: the composition $F^pK^\bullet \to J^{p, \bullet}$
is homotopic to zero because $J^{p, \bullet}$ is acyclic and K-injective
(Derived Categories, Lemma \ref{derived-lemma-K-injective}).
Since the objects $J^{p, n - 1}$ are injective and since
$F^pK^n \to K^n \to I^n$ are injective morphisms, we
can lift the maps $F^pK^n \to J^{p, n - 1}$ giving the homotopy
to a map $h^n : I^n \to J^{p, n - 1}$. Then we set $s^p$
equal to $h \circ \text{d} + \text{d} \circ h$.
(Warning: It will not be the case that $t^p = s^p \circ \alpha^p$,
so we have to be careful not to use this below.)

\medskip\noindent
Consider
$$
J^\bullet = I^\bullet \times \prod\nolimits_p J^{p, \bullet}
$$
Because products in $D(\mathcal{A})$ are given by taking
products of K-injective complexes
(Lemma \ref{lemma-derived-products})
and since $J^{p, \bullet}$
is isomorphic to $0$ in $D(\mathcal{A})$ we see that
$J^\bullet \to I^\bullet$ is an isomorphism in $D(\mathcal{A})$.
Consider the map
$$
j = i \times (s^p \circ i)_{p \in \mathbf{Z}} :
K^\bullet
\longrightarrow
I^\bullet \times \prod\nolimits_p J^{p, \bullet} = J^\bullet
$$
By our remarks above this is a quasi-isomorphism. It is also injective.
For $p \in \mathbf{Z}$ we let $F^pJ^\bullet \subset J^\bullet$ be
$$
\Im\left(
\alpha^p \times (t^{p'} \circ \alpha^{pp'})_{p' \leq p} :
I^{p, \bullet}
\to
I^\bullet \times \prod\nolimits_{p' \leq p} J^{p', \bullet}
\right)
\times \prod\nolimits_{p' > p} J^{p', \bullet}
$$
This complex is isomorphic to the complex
$I^{p, \bullet} \times \prod_{p' > p} J^{p, \bullet}$
as $\alpha^{pp} = \text{id}$ and $t^p$ is injective.
Hence $F^pJ^\bullet$ is quasi-isomorphic to $I^{p, \bullet}$ (argue
as above). We have $j(F^pK^\bullet) \subset F^pJ^\bullet$ because
of the commutativity of the diagram above. The corresponding
map of complexes $F^pK^\bullet \to F^pJ^\bullet$ is a quasi-isomorphism
by what we just said. Finally, to see that
$F^{p + 1}J^\bullet \subset F^pJ^\bullet$
use that $\alpha^{p + 1p} \circ \alpha^{pp'} = \alpha^{p + 1p'}$
and the commutativity of the first displayed diagram
in the first paragraph of the proof.

\medskip\noindent
We claim that $j : K^\bullet \to J^\bullet$ is a solution to the
problem posed by the lemma. Namely, $F^pJ^n$ is an injective object
of $\mathcal{A}$ because it is isomorphic to
$I^{p, n} \times \prod_{p' > p} J^{p', n}$ and products of
injectives are injective. Then the injective map $F^pJ^n \to J^n$
splits and hence the quotient $J^n/F^pJ^n$ is injective as well
as a direct summand of the injective object $J^n$.
Similarly for $F^pJ^n/F^{p'}J^n$. This in particular means
that $0 \to F^pJ^\bullet \to J^\bullet \to J^\bullet/F^pJ^\bullet \to 0$
is a termwise split short exact sequence of complexes, hence defines
a distinguished triangle in $K(\mathcal{A})$ by fiat.
Since $J^\bullet$ and $F^pJ^\bullet$ are K-injective complexes
we see that the same is true for $J^\bullet/F^pJ^\bullet$
by Derived Categories, Lemma \ref{derived-lemma-triangle-K-injective}.
A similar argument shows that $F^pJ^\bullet/F^{p'}J^\bullet$
is K-injective. By construction $j : K^\bullet \to J^\bullet$
and the induced maps $F^pK^\bullet \to F^pJ^\bullet$ are
quasi-isomorphisms. Using the long exact cohomology sequences
of the complexes in play we find that the same holds for
$K^\bullet/F^pK^\bullet \to J^\bullet/F^pJ^\bullet$ and
$F^pK^\bullet/F^{p'}K^\bullet \to F^pJ^\bullet/F^{p'}J^\bullet$.
\end{proof}

\bibliographystyle{plain}
\bibliography{my}
\end{document}
